\documentclass[11pt]{amsart}
\usepackage{amsmath,amssymb}
\usepackage{tikz}
\usepackage{hyperref}
\usetikzlibrary{matrix,arrows}
\newtheorem{theorem}{Theorem}
\newtheorem{lemma}{Lemma}
\newtheorem{proposition}{Proposition}
\newtheorem{corollary}{Corollary}
\theoremstyle{definition}
\newtheorem{definition}{Definition}
\newcommand{\problemhead}[1]{\medskip\noindent\textbf{Problem #1.} }
\newcommand{\exercisehead}[1]{\medskip\noindent\textbf{Exercise #1.} }
\newcommand{\solutionhead}[1]{\noindent\textbf{Solution #1.} }
\title{Introduction to Smooth Manifolds: The Cotangent Bundle Solutions}
\begin{document}
\maketitle
\section{The Cotangent Bundle}

\subsection*{Covectors}

\exercisehead{11.2} For linear $\omega : V \to \mathbb{R}$, $\omega \in V^*$
\[
\begin{gathered}
\omega(E_j) = \omega( \delta^i_{ \, j} E_i) = \delta^i_{ \, j } \omega(E_i) = \omega(E_i) \epsilon^i(E_j) \\
\Longrightarrow \omega = \omega(E_i)\epsilon^i
\end{gathered}
\]
So $\omega$ spanned by $\epsilon^i$, \, $i = 1 \dots n$

Suppose $0 = \omega_i \epsilon^i$.  Then $\forall \, x \in V$, \[
\omega_i \epsilon^i(x^j E_j) = \omega_i x^j \delta^i_{ \, j } = \omega_i x^i = \omega(E_i) x^i = \omega(x) = 0
\]
Then $\omega=0$ since $\omega(x) =0, \, \forall\, x \in V$.  So $\epsilon^i$ form a linear basis of $V^*$.

By theorem, $\text{dim}{ V^*} = \text{dim}{V}$



\end{document}
